% Vietnamese translation for OI Public Library
% Source-PDF: IOI中国国家候选队论文/国家集训队2020论文集.pdf
\documentclass[11pt]{article}
\usepackage{oipl-vn}

\title{Tập luận văn đội tuyển quốc gia Trung Quốc năm 2020}
\author{Tập luận văn đội tuyển quốc gia Olympic Tin học}
\date{China Computer Federation, 2020}

\begin{document}
\maketitle

\origtitle{Tập luận văn đội tuyển quốc gia Trung Quốc Olympic Tin học năm 2020}
\sourcepath{IOI China candidate team papers / National training team 2020 collection.pdf}

\renewcommand{\abstractname}{Tóm tắt}
\begin{abstract}
PDF nguồn là tập hợp luận văn đội tuyển quốc gia Trung Quốc năm 2020. Tệp được
tạo bằng XeTeX và có lớp văn bản Unicode đọc được: mục lục, phần mở đầu và các
trang thân bài đã kiểm tra đều trích xuất tốt bằng \texttt{pdftotext}. Bản
LaTeX này chuyển ngữ phần nhận diện tập sách và toàn bộ mục lục, đồng thời ghi
lại các chủ đề chính để phục vụ bản dịch chi tiết hơn.
\end{abstract}

\section{Thông tin đầu sách}

Tập sách có tiêu đề \emph{Tập luận văn đội tuyển quốc gia Trung Quốc Olympic
Tin học năm 2020}. Tệp PDF gồm 196 trang và được tạo bằng XeTeX. Trang bìa
không ghi riêng tên huấn luyện viên trong phần văn bản trích xuất được.

\section{Mục lục đã dịch}

\begin{center}
\small
\begin{tabular}{r p{0.52\linewidth} p{0.24\linewidth}}
\textbf{Trang} & \textbf{Bài viết} & \textbf{Tác giả} \\
\hline
3 & Bàn về tối ưu hóa bài toán chia khối cổ điển bằng fractional cascading & Jiang Mingrun \\
12 & Bàn về cây thống trị và ứng dụng & Chen Sunli \\
22 & Báo cáo ra đề \emph{Ghép số} & Jiang Xunchi \\
32 & Báo cáo ra đề \emph{Khối liên thông nhỏ nhất} & Pan Junyue \\
42 & Giới thiệu đơn giản về nguyên lý chuyển vị & Chen Yu \\
52 & Bàn về duy trì động giá trị cực trị của hàm số & Li Baitian \\
63 & Báo cáo ra đề \emph{Trò chơi trên đồ thị} & Zuo Junchi \\
74 & Một loại DFT không bình thường & Zhou Yuyang \\
87 & Bài toán rừng hướng vào nhỏ nhất & Zhang Zheyu \\
102 & Bàn về suffix automaton nén & Xu Yixuan \\
117 & Bàn về nimber và thuật toán đa thức & Luo Yuxiang \\
133 & Các bài toán liên quan tới định thức trong thi tin học & Wang Zhanpeng \\
151 & Thuật toán hiệu quả để tính diện tích vùng ghép trên mặt cầu & Zhou Renfei \\
173 & Bàn về ứng dụng đơn giản của lý thuyết nhóm trong thi tin học & Yu Haoxiang
\end{tabular}
\end{center}

\section{Chủ đề chính}

Từ mục lục và phần đầu của bài thứ nhất, tập năm 2020 bao phủ các nhóm chủ đề
sau:
\begin{itemize}
  \item fractional cascading và tối ưu bài toán chia khối;
  \item cây thống trị, khối liên thông nhỏ nhất và trò chơi trên đồ thị;
  \item nguyên lý chuyển vị, DFT, định thức và thuật toán đa thức;
  \item duy trì động cực trị hàm số và rừng hướng vào nhỏ nhất;
  \item suffix automaton nén, nimber và lý thuyết nhóm;
  \item hình học trên mặt cầu và tính diện tích vùng ghép.
\end{itemize}

\section{Ghi chú về trích xuất}

Phần thân của tập 2020 có thể trích xuất được mà không cần OCR. Bài đầu tiên
giới thiệu fractional cascading như một kỹ thuật tăng tốc tìm kiếm nhị phân
trên nhiều dãy có tổ chức, rồi dùng nó để tối ưu một bài toán chia khối cổ
điển. Điều này cho thấy tập 2020 phù hợp để mở rộng thành bản dịch từng bài,
bao gồm công thức, hình và bảng, bằng cách xử lý trực tiếp văn bản Unicode
trích xuất được.

\end{document}
